\documentclass[11pt]{article}

\usepackage[margin=1in]{geometry}
\usepackage{amsmath,amssymb,amsthm}
\usepackage{booktabs}
\usepackage{hyperref}
\usepackage{cleveref}

\newtheorem{theorem}{Theorem}
\newtheorem{corollary}[theorem]{Corollary}
\newtheorem{lemma}[theorem]{Lemma}
\newtheorem{proposition}[theorem]{Proposition}
\theoremstyle{definition}
\newtheorem{definition}[theorem]{Definition}
\theoremstyle{remark}
\newtheorem{remark}[theorem]{Remark}

\DeclareMathOperator*{\argmin}{arg\,min}
\newcommand{\bs}{\mathbf{s}}
\newcommand{\bD}{\mathbf{D}}
\newcommand{\bx}{\mathbf{x}}
\newcommand{\NML}{\operatorname{NML}}
\newcommand{\NLL}{\operatorname{NLL}}
\newcommand{\RL}{\operatorname{RL}}
\newcommand{\Hb}{H_b}

\title{Model Selection Stabilization for Relative-Periodic Decomposition\\of Cellular Automaton Spacetimes}
\author{}
\date{}

\begin{document}
\maketitle

\begin{abstract}
We prove that model selection over relative-periodic backgrounds of cellular automaton (CA) spacetimes stabilizes to a unique period as the observation window grows.
The argument rests on a dimension-agnostic \emph{orbit-class reduction}: fitting a relative-periodic background to a binary spacetime decomposes into independent majority voting on orbit classes (Theorem~\ref{thm:projection}), partition refinement along constant-velocity chains makes higher periods monotonically more expressive (Theorem~\ref{thm:monotonicity}), and a Bernoulli NML criterion with $O(\log T)$ complexity penalty against $\Theta(T)$ data-fit cost yields eventual stabilization of the selected model (Theorem~\ref{thm:stabilization}).
We also prove that background period recovery is impossible in general without additional assumptions (Theorem~\ref{thm:nonident}).
The theory applies identically to 1D, 2D, and 3D automata.

Cross-dimensional experiments confirm the stabilization: selected periods lock for 1D elementary CA (rules 30, 54, 110), 2D totalistic rules (773 non-trivial from 1{,}050 candidates), and 3D totalistic rules---with margins growing after stabilization.
As an application, the framework identifies 2D rules with persistent structured projection residuals, including one (S37/B11) exhibiting extensive residual scaling verified at 400 steps, across multiple seeds, and at grid sizes up to $192 \times 192$.

\medskip\noindent
\textbf{Keywords:} cellular automata, normalized maximum likelihood, model selection stabilization, symmetry projection, projection residual decomposition
\end{abstract}

%======================================================================
\section{Introduction}
%======================================================================

Cellular automata (CA) generate complex spatiotemporal patterns from simple local rules.
Decomposing these patterns into regular and irregular components is a central problem in CA theory, with applications to understanding emergent computation, classifying rule complexity, and identifying coherent structures.

The computational mechanics program \cite{crutchfield1997,hanson1992} builds local epsilon-machine models that identify domains, particles, and their interactions.
Redeker~\cite{redeker2017} formalized particle catalogs via de Bruijn diagrams.
Rupe and Crutchfield~\cite{rupe2018} generalized local causal states to arbitrary dimensions.
Shalizi et al.~\cite{shalizi2006} developed automatic coherent-structure filters.
These methods build rich local models at significant computational cost.

We take a different approach: rather than building local models, we ask a global model-selection question.
\emph{Given a finite family of relative-periodic templates, which template best describes the data?}
We show that this question has a clean answer---the NML-optimal period stabilizes---and that the stabilization proof is dimension-agnostic, applying uniformly to 1D, 2D, and 3D automata.

\subsection{Contributions}

\begin{enumerate}
    \item \textbf{Orbit-class reduction} (Theorem~\ref{thm:projection}): fitting a relative-periodic background reduces to independent majority voting on orbit classes, yielding the unique Hamming-optimal projection in $O(n)$ time.
    \item \textbf{Monotonicity under velocity matching} (Theorem~\ref{thm:monotonicity}): higher periods achieve lower defect rates along constant-velocity divisibility chains, necessitating complexity control.
    \item \textbf{Model selection stabilization} (Theorem~\ref{thm:stabilization}): under ergodic defect rates, the NML-selected period stabilizes with growing margins---proved for any spatial dimension, with no assumption on defect geometry.
    \item \textbf{Nonidentifiability} (Theorem~\ref{thm:nonident}): background period recovery is impossible in general; periodic defect structure can be absorbed by higher-period templates.
    \item \textbf{Cross-dimensional experiments}: empirical stabilization observed for 1D (ECA 30/54/110), 2D (773 totalistic rules), and 3D (totalistic rules on a 3-torus).
    \item \textbf{Application}: identification of persistent structured projection residuals in 2D rules, with extensive scaling in S37/B11.
\end{enumerate}

\subsection{Relationship to Prior Work}

\begin{table}[h]
\centering
\begin{tabular}{lll}
\toprule
Aspect & Computational Mechanics \cite{crutchfield1997,hanson1992,rupe2018} & This Work \\
\midrule
Model class & Local (epsilon-machine) & Global (relative-periodic) \\
Selection criterion & Statistical complexity & Bernoulli NML \\
Stabilization proof & --- & Theorem~\ref{thm:stabilization} \\
Nonidentifiability & --- & Theorem~\ref{thm:nonident} \\
Dimensionality & Case-by-case & Dimension-agnostic \\
Cost & Higher & $O(n)$ per model \\
\bottomrule
\end{tabular}
\end{table}

Packard and Wolfram~\cite{packard1985} surveyed 2D rules qualitatively.
Boccara and Roger~\cite{boccara1999} identified period-2 families.
Zenil~\cite{zenil2010} used compression for CA classification.
Our contribution is a \emph{stabilization theorem} for period selection that applies across dimensions, with the surveys as empirical evidence.

%======================================================================
\section{Theory}
%======================================================================

\subsection{Relative-Periodic Model Class}

\begin{definition}[Relative-periodic model class]\label{def:model}
Given a binary spacetime $U \in \{0,1\}^{T \times D_1 \times \cdots \times D_n}$, shift vector $\bs = (s_1, \ldots, s_n) \in \mathbb{Z}^n$, and period $p \geq 1$, the \emph{relative-periodic model class} $\mathcal{B}(p, \bs)$ consists of all fields $B$ satisfying
\[
B[t + p,\; (x_1 + s_1) \bmod D_1,\; \ldots,\; (x_n + s_n) \bmod D_n] = B[t, x_1, \ldots, x_n]
\]
for all valid indices.
This constraint partitions spacetime into \emph{orbit classes} $\{O_j\}_{j=1}^{k}$ where $k = p \prod_i D_i$.
Within each class, $B$ is constant.
The definition is identical for $n = 1, 2, 3$, or any dimension.
\end{definition}

\subsection{Orbit-Class Reduction}

\begin{theorem}[Optimal Hamming Projection]\label{thm:projection}
Let $n_j^{(1)} = |\{(t,\bx) \in O_j : U[t,\bx] = 1\}|$ and $n_j^{(0)} = |O_j| - n_j^{(1)}$.
The Hamming distance $d(U, B) = |\{(t,\bx) : U[t,\bx] \neq B[t,\bx]\}|$ is minimized over $\mathcal{B}(p, \bs)$ by setting
\[
b_j = \begin{cases}
1 & \text{if } n_j^{(1)} > n_j^{(0)}, \\
0 & \text{if } n_j^{(0)} > n_j^{(1)}, \\
\text{either} & \text{if } n_j^{(0)} = n_j^{(1)}.
\end{cases}
\]
The minimum defect count is $\sum_j \min(n_j^{(0)}, n_j^{(1)})$.
The minimizer is unique if and only if no orbit class has $n_j^{(0)} = n_j^{(1)}$.
\end{theorem}

\begin{proof}
The Hamming distance decomposes additively over orbit classes: $d(U,B) = \sum_j d_j$ where $d_j$ counts disagreements within $O_j$.
Since $B$ is constant on each $O_j$, choosing $b_j$ independently minimizes each $d_j = \min(n_j^{(0)}, n_j^{(1)})$.
\end{proof}

This is the \emph{orbit-class reduction}: the spatiotemporal fitting problem decomposes into $k$ independent binary estimation problems.
The reduction is exact and holds in any dimension.
The optimal projection is computable in $O(|U|)$ time via a single pass over the data.

\begin{definition}[Defect mask and rate]\label{def:defect}
The \emph{defect mask} is $M = U \oplus B^*$ where $B^*$ is the optimal projection.
The \emph{defect rate} is $r = \|M\|_1 / |U|$.
\end{definition}

\subsection{Monotonicity and Overcapacity}

\begin{theorem}[Monotonicity under Velocity-Matched Refinement]\label{thm:monotonicity}
Let $(p_1, \bs_1)$ and $(p_2, \bs_2)$ be two relative-periodic models.
If there exists an integer $m \geq 1$ such that $p_2 = m \cdot p_1$ and $s_2^{(i)} \equiv m \cdot s_1^{(i)} \pmod{D_i}$ for each spatial dimension $i$, then the orbit partition under $(p_2, \bs_2)$ refines the partition under $(p_1, \bs_1)$, and
\[
d^*(p_2, \bs_2) \leq d^*(p_1, \bs_1).
\]
The condition $s_2 = m \cdot s_1 \pmod{D}$ means the two models share the same \emph{characteristic velocity} $\bs/p$.
In particular, for shift $\bs = \mathbf{0}$, monotonicity holds along all divisibility chains $p, 2p, 3p, \ldots$
\end{theorem}

\begin{proof}
Under the stated condition, the periodicity map $\tau_2\colon (t, \bx) \mapsto (t + p_2, \bx + \bs_2 \bmod \bD)$ equals $\tau_1^m$, the $m$-fold composition of the $(p_1, \bs_1)$ map.
Therefore, the $\tau_2$-orbit of any point is a subset of its $\tau_1$-orbit, so the $\tau_2$-partition refines the $\tau_1$-partition.
By Theorem~\ref{thm:projection}, optimizing over a finer partition can only reduce total disagreement.
\end{proof}

\begin{remark}
The velocity-matching condition is necessary in general.
For example, on a ring of width~4, the models $(p{=}1, s{=}1)$ and $(p{=}2, s{=}1)$ share the same shift but \emph{not} the same velocity ($1/1 \neq 1/2$).
The checkerboard spacetime $U[t,x] = (t+x) \bmod 2$ has $d^*(1,1) = 0$ but $d^*(2,1) = 4$, violating monotonicity.
\end{remark}

\begin{corollary}[Overcapacity]
Along any constant-velocity divisibility chain, the residual count is monotonically non-increasing.
This creates an \emph{overcapacity problem}: naive residual-count minimization always prefers the highest available period.
This motivates the complexity-penalized criterion below.
\end{corollary}

\subsection{Bernoulli NML Criterion}

The orbit-class reduction (Theorem~\ref{thm:projection}) shows that each model $(p, \bs)$ decomposes the data into $k = p \prod_i D_i$ independent Bernoulli estimation problems.
This motivates a \emph{Bernoulli NML} criterion where the statistical model family is exactly the one implied by the orbit-class structure.

\begin{definition}[Orbit-class Bernoulli model]\label{def:bernoulli}
For model $(p, \bs)$ with $k$ orbit classes, each orbit class $O_j$ has $n_j = |O_j|$ observations and empirical frequency $\hat{\theta}_j = n_j^{(1)} / n_j$.
The Bernoulli maximum-likelihood estimate for the $j$-th class is $\hat{\theta}_j$, with negative log-likelihood
\[
\NLL(p, \bs) = \sum_{j=1}^{k} n_j \, \Hb(\hat{\theta}_j),
\]
where $\Hb(\theta) = -\theta \log_2 \theta - (1-\theta) \log_2(1-\theta)$ is binary entropy (with $0 \log 0 = 0$).
\end{definition}

\begin{definition}[Bernoulli NML score]\label{def:nml}
The normalized maximum likelihood score for $k$ independent Bernoulli classes is
\[
\NML(p, \bs) = \underbrace{\NLL(p, \bs)}_{\text{data fit}} + \underbrace{\sum_{j=1}^{k} \tfrac{1}{2} \log_2 n_j}_{\text{asymptotic parametric complexity}}.
\]
The complexity term $\frac{1}{2} \log_2 n_j$ is the asymptotic NML parametric complexity for class $j$ with $n_j$ observations \cite{grunwald2007}.
The combined score is an asymptotic approximation to the normalized maximum likelihood code for the orbit-class Bernoulli model.
It is not exact NML (the exact normalizing constants involve sums of binomial coefficients), but the $O(k)$ additive error does not affect asymptotic model selection since the NLL term is $\Theta(T)$.
\end{definition}

\begin{remark}[Relationship to Hamming projection]
Majority vote (Theorem~\ref{thm:projection}) minimizes Hamming distance; the Bernoulli MLE $\hat{\theta}_j$ minimizes negative log-likelihood.
These serve different purposes: majority vote produces the optimal background decomposition, while NML selects among competing decompositions.
The two coincide when $\hat{\theta}_j \in \{0, 1\}$ (pure orbit classes) and diverge most when $\hat{\theta}_j \approx 1/2$ (maximally noisy classes).
\end{remark}

\begin{remark}[No geometry dependence]
Unlike run-length or LZ4 codelengths, the NML score depends only on the orbit-class statistics $(n_j, n_j^{(1)})$---it is independent of the spatial arrangement of defects and of the traversal order.
This makes the criterion intrinsic to the orbit-class structure.
\end{remark}

\subsection{Stabilization Theorem}

This is the main theoretical result.

\begin{theorem}[Stabilization of Bernoulli NML Selection]\label{thm:stabilization}
Let $\mathcal{C} = \{(p_1, \bs_1), \ldots, (p_m, \bs_m)\}$ be a fixed finite candidate set.
For each candidate $c \in \mathcal{C}$, let $k_c = p_c \prod_i D_i$ be the number of orbit classes and let $n_j(T)$ denote the size of the $j$-th orbit class at observation length $T$.
Assume:
\begin{itemize}
    \item[\textup{(A1)}] \emph{Ergodic orbit-class frequencies}: for each candidate $c$ and each orbit class $j$, the empirical frequency $\hat{\theta}_j(T) \to \theta_j^*$ as $T \to \infty$.
\end{itemize}
Then:
\begin{enumerate}
    \item[\textup{(i)}] The per-site NLL rate converges: $\NLL(c, T) / T \to \lambda_c$ where $\lambda_c = \lim_{T\to\infty} \frac{1}{T}\sum_j n_j(T) \Hb(\theta_j^*)$.
    \item[\textup{(ii)}] The parametric complexity is $O(\log T)$: specifically, $\sum_j \frac{1}{2}\log_2 n_j(T) = \frac{k_c}{2} \log_2 \frac{T}{p_c} + O(1)$ since each orbit class has $n_j = \lfloor T/p_c \rfloor$ or $\lceil T/p_c \rceil$ observations.
    \item[\textup{(iii)}] The NML-selected candidate $c^*(T) = \argmin_{c \in \mathcal{C}} \NML(c, T)$ stabilizes: there exists $T_0$ such that $c^*(T) = c^*(T_0)$ for all $T > T_0$.
    \item[\textup{(iv)}] The stabilized selection minimizes the asymptotic per-site NLL rate $\lambda_c$, breaking ties by parametric complexity (fewer parameters preferred).
\end{enumerate}
\end{theorem}

\begin{proof}
Under (A1), the NML score for candidate $c$ decomposes as
\[
\NML(c, T) = \NLL(c, T) + \sum_j \tfrac{1}{2}\log_2 n_j(T) = \lambda_c \cdot T + \frac{k_c}{2}\log_2 T + O(1).
\]
For any two candidates $c_1, c_2$:
\[
\NML(c_1, T) - \NML(c_2, T) = (\lambda_{c_1} - \lambda_{c_2}) \cdot T + \frac{k_{c_1} - k_{c_2}}{2} \log_2 T + O(1).
\]
\textbf{Case 1:} $\lambda_{c_1} \neq \lambda_{c_2}$.
The linear term dominates for large $T$, so the candidate with smaller $\lambda$ wins permanently.

\textbf{Case 2:} $\lambda_{c_1} = \lambda_{c_2}$.
The logarithmic term dominates, and the simpler model (fewer parameters) wins permanently.

Since $|\mathcal{C}|$ is finite, all pairwise comparisons stabilize, giving a unique winner for $T > T_0$.
\end{proof}

\begin{corollary}[Dimension-agnostic]
Theorem~\ref{thm:stabilization} makes no reference to spatial dimension $n$.
It applies identically to 1D rings, 2D tori, 3D tori, or any periodic lattice.
\end{corollary}

\begin{remark}[Run-length encoding alternative]\label{rem:rl}
An alternative MDL criterion using $\frac{k}{2} \log_2(T/p) + L_{\RL}(M^*)$ (BIC-type penalty plus run-length defect encoding) satisfies an analogous stabilization theorem, provided the per-site RL rate $\ell_c = \lim L_{\RL} / N(T)$ converges.
For finite deterministic CA, this convergence can be proved (see Appendix~\ref{app:rl}).
We use the NML score as the primary selection criterion because it converges \emph{unconditionally} (no ties assumption needed), and report RL codelength as a secondary geometric diagnostic (Proposition~\ref{prop:rl-separation}).
\end{remark}

\subsection{Nonidentifiability}

\begin{theorem}[Nonidentifiability of Background Period]\label{thm:nonident}
For any period $p_0$, shift $\bs$, and MDL-type criterion with $O(\log T)$ penalty and $\Theta(T)$ data-fit cost, there exist spacetimes for which the criterion does not select $p_0$ for any $T$ beyond some threshold.
Specifically, if the projection residual under $(p_0, \bs)$ has its own periodic structure with period $k > 1$, then the candidate $(k \cdot p_0, k \cdot \bs \bmod \bD)$ absorbs the residual into its template and eventually achieves a lower total score.
\end{theorem}

The observed spacetime does not intrinsically separate into ``background'' plus ``residual'' without additional assumptions.
The selected period is the one that best compresses the \emph{entire} spacetime, not necessarily the one matching an external notion of ``true background period.''

\subsection{Identifiability under Aperiodic Residuals}

\begin{definition}[True background period]\label{def:truep}
The \emph{true background period} $p_0$ of a CA spacetime is the smallest period such that the projection residual $M^*_{p_0}$ has vanishing structure: all orbit classes under $p_0$ have asymptotic frequency $\theta_j^* \in \{0, 1\}$.
\end{definition}

\begin{theorem}[Identifiability under Aperiodic Residuals]\label{thm:ident}
Let $p_0$ be the true background period and suppose residuals are aperiodic in the following sense: for all orbit classes $j$ under model $p_0$, the empirical frequency $\theta_j^* \in \{0, 1\}$, but for any strict multiple $p = m \cdot p_0$ ($m > 1$), the refinement of some orbit class produces sub-classes with the same frequency $\theta_j^* \in \{0, 1\}$.
Then:
\begin{enumerate}
    \item[\textup{(i)}] $\NLL(p, T) = \NLL(p_0, T) + o(T)$ for any multiple $p = m \cdot p_0$, because splitting a pure orbit class ($\theta^* \in \{0,1\}$) into sub-classes does not reduce NLL.
    \item[\textup{(ii)}] The parametric complexity for $p$ exceeds that of $p_0$: $\frac{k_p}{2}\log_2 T > \frac{k_{p_0}}{2}\log_2 T$ since $k_p = m \cdot k_{p_0}$.
    \item[\textup{(iii)}] Therefore NML selects $p_0$ over all its strict multiples for large $T$.
\end{enumerate}
\end{theorem}

\begin{proof}
When $\theta_j^* \in \{0,1\}$, we have $\Hb(\theta_j^*) = 0$, so the NLL contribution from class $j$ is zero.
Splitting $O_j$ into $m$ sub-classes $O_{j_1}, \ldots, O_{j_m}$ with the same frequency preserves $\Hb = 0$ for each sub-class, so NLL is unchanged.
But the complexity increases by $\frac{m-1}{2}\log_2(T/p)$ per split class.
Since $k_p > k_{p_0}$, the complexity difference $\frac{k_p - k_{p_0}}{2}\log_2 T \to +\infty$, and NML penalizes the larger model.
\end{proof}

\begin{remark}
Theorems~\ref{thm:nonident} and~\ref{thm:ident} are complementary: when residuals are aperiodic, NML recovers the true period (Theorem~\ref{thm:ident}); when residuals have their own periodic structure, NML absorbs them into higher-period templates (Theorem~\ref{thm:nonident}).
This is \emph{principled} behavior: the data genuinely admits a simpler (zero-residual) description at the higher period.
The distinction between ``background periodicity'' and ``residual periodicity'' is not intrinsic to the observation.
\end{remark}

\subsection{Geometric Diagnostics: Run-Length and LZ4}

The NML criterion (Definition~\ref{def:nml}) is the primary model selector.
For geometric analysis of defect masks \emph{after} model selection, we use two compression-based diagnostics:

\textbf{Run-length codelength} $L_{\RL}(M)$: Elias-gamma code over run lengths of the flattened defect mask.
This measures spatial clustering: a contiguous defect block costs $O(\log n)$ bits, while $w$ scattered defects cost $\Omega(w \log(n/w))$ bits.

\textbf{LZ4 codelength} $L_{\mathrm{LZ4}}(M)$: LZ4 compression of the packed defect mask.
A practical compressor proxy.

\begin{proposition}[Run-Length Separation]\label{prop:rl-separation}
Two binary masks of the same size and Hamming weight can have run-length codelengths differing by a factor of $\Omega(\log n)$.
\end{proposition}

\begin{proof}[Proof sketch]
A mask of weight $w$ with a single contiguous block has RL codelength $O(\log n)$.
A mask of the same weight with $w$ isolated defects has RL codelength $\Omega(w \log(n/w))$.
\end{proof}

These diagnostics are \emph{not} used for model selection---they depend on traversal order and do not correspond to the NML data-fit term.
They are useful for characterizing defect geometry after the NML-selected model is fixed.

%======================================================================
\section{Cross-Dimensional Experiments}
%======================================================================

We test Theorem~\ref{thm:stabilization} by sweeping the observation length $T$ and tracking the NML-selected period across 1D, 2D, and 3D cellular automata.

\subsection{1D Elementary CA}

ECA rules 30, 54, and 110 on a ring of width 192, scanning shifts $\pm 6$, periods 1--10.

\begin{table}[h]
\centering
\caption{NML-optimal decomposition ($T = 144$, shift $= 0$).}
\begin{tabular}{lccccc}
\toprule
Rule & Best $(s, p)$ & Defect Rate & NLL Bits & Complexity & NML Bits \\
\midrule
54  & $(0, 4)$ & 20.1\% & 17{,}477 & 1{,}985 & 19{,}462 \\
110 & $(0, 7)$ & 30.8\% & 22{,}842 & 2{,}931 & 25{,}774 \\
30  & $(5, 1)$ & 46.3\% & 27{,}489 & 688     & 28{,}177 \\
\bottomrule
\end{tabular}
\end{table}

The known complexity hierarchy Rule~54 $<$ Rule~110 $<$ Rule~30 \cite{crutchfield1997,hanson1992,redeker2017} is recovered under NML ranking.

\begin{table}[h]
\centering
\caption{Period stabilization (1D).}
\begin{tabular}{lccccccc}
\toprule
Rule & $T{=}50$ & $T{=}100$ & $T{=}200$ & $T{=}400$ & $T{=}600$ & $T{=}800$ & Margin at $T{=}800$ \\
\midrule
ECA-54  & 4 & 4 & 4 & 4 & 4 & 4 & 1{,}993 \\
ECA-110 & 7 & 7 & 7 & 7 & 7 & 7 & 18{,}378 \\
ECA-30  & 1 & 1 & 1 & 1 & 1 & 1 & 598 \\
\bottomrule
\end{tabular}
\end{table}

All three 1D rules show \emph{immediate} stabilization---the NML-selected period is constant from $T{=}50$ onward, with positive margins at every point.
Margins grow monotonically, consistent with Theorem~\ref{thm:stabilization}.
Rule~110 shows the strongest separation (margin 18{,}378 bits at $T{=}800$), consistent with its pronounced period-7 Ether background.

\subsection{2D Totalistic Rules}

\subsubsection{Survey}

We survey 1{,}050 range-threshold totalistic rules on a $48 \times 48$ torus (Moore neighborhood, 40 steps, density 0.5, seed 11).
A cell survives if $s_{\text{lo}} \leq n \leq s_{\text{hi}}$ and is born if $b_{\text{lo}} \leq n \leq b_{\text{hi}}$, with range widths $\leq 4$.
Of 1{,}050 candidates, 773 produce non-trivial dynamics (end density between 2\% and 98\%).
Of these, 172 have nonzero best-fit shift (drifting periodic structure).

Top rules by NML score are near-fixed-point patterns with ${<}1\%$ defect rates---NML correctly identifies these as the simplest descriptions.
Period-2 oscillators (159 rules) appear at NML rank ${\sim}20$, with no period ${\geq}3$ selected at this short horizon (40 steps).
Period distribution: 614 period-1, 159 period-2.

\subsubsection{Multi-Seed Robustness}

We verified 6 rules across 10 seeds ($64 \times 64$, 60 steps):

\begin{table}[h]
\centering
\caption{Multi-seed robustness.}
\begin{tabular}{lcccc}
\toprule
Rule & Modal Period & Mean Rate & CV & Period Consistent \\
\midrule
S14/B11 & 1 & 0.59\% & 8.4\%  & 10/10 \\
S25/B12 & 2 & 0.58\% & 15.6\% & 8/10 \\
S66/B36 & 2 & 0.63\% & 10.5\% & 10/10 \\
S24/B11 & 2 & 2.78\% & 9.4\%  & 10/10 \\
S11/B37 & 4 & 2.75\% & 18.7\% & 7/10 \\
S37/B11 & 2 & 4.41\% & 12.4\% & 10/10 \\
\bottomrule
\end{tabular}
\end{table}

Rules with near-zero NML margins (S25/B12, S11/B37) show seed-dependent period selection at short horizons---consistent with Theorem~\ref{thm:stabilization}, which requires the linear NLL terms to dominate.

\subsubsection{Period Stabilization (2D)}

\begin{table}[h]
\centering
\caption{Period stabilization (2D).}
\begin{tabular}{lccccccc}
\toprule
Rule & $T{=}60$ & $T{=}100$ & $T{=}200$ & $T{=}400$ & $T{=}600$ & $T{=}800$ & Margin at $T{=}800$ \\
\midrule
S24/B11 & 1 & 1 & 1 & 2 & 2 & 2 & 13{,}115 \\
S11/B37 & 2 & 2 & 2 & 4 & 4 & 4 & 22{,}352 \\
S37/B11 & 1 & 1 & 2 & 2 & 2 & 2 & 16{,}006 \\
\bottomrule
\end{tabular}
\end{table}

The 2D persistent-residual rules show \emph{progressive} stabilization: the NML-selected period increases in discrete jumps ($1 \to 2$ for S24/B11, $2 \to 4$ for S11/B37, $1 \to 2$ for S37/B11), then locks with growing margins.
This is the regime described by Theorem~\ref{thm:nonident} (nonidentifiability): periodic residual structure is absorbed into higher-period templates, and NML correctly identifies the period that best compresses the entire spacetime.

Notably, NML selects \emph{lower} periods than the previous MDL formulation (which used RL codelength as data fit).
For example, S24/B11 stabilizes to period~2 under NML vs period~6 under MDL---the RL artifact inflated the data-fit term for higher periods, making them appear artificially competitive.

\subsection{3D Totalistic Rules}

Diamoeba3d rule (survive 5--8, birth 5--8) on a $16 \times 16 \times 16$ 3-torus, shift $(0,0,0)$:

\begin{table}[h]
\centering
\caption{Period stabilization (3D).}
\begin{tabular}{cccc}
\toprule
$T$ & Period & NML Bits & Margin \\
\midrule
10 & 6 & 26{,}406  & 5{,}560 \\
20 & 6 & 75{,}703  & 3{,}055 \\
40 & 1 & 158{,}154 & 3{,}505 \\
60 & 1 & 232{,}856 & 4{,}872 \\
80 & 1 & 303{,}209 & 5{,}840 \\
\bottomrule
\end{tabular}
\end{table}

The initial $T{=}10$--$20$ selection of period~6 is a small-sample artifact (10--20 time steps with a period-6 model means only 1--3 full cycles per orbit class, where overfitting is expected).
By $T{=}40$, NML stabilizes to period~1 with margins growing monotonically---consistent with Theorem~\ref{thm:stabilization} in 3D.
The 3D-life rule (survive 4--5, birth 5) dies out, providing a null result.

\subsection{Summary of Stabilization Evidence}

\begin{table}[h]
\centering
\caption{Cross-dimensional stabilization summary.}
\begin{tabular}{llll}
\toprule
Dimension & Rules Tested & Stabilization Onset & Behavior \\
\midrule
1D & ECA 30, 54, 110            & Immediate ($T{=}50$) & Period locked from first measurement \\
2D & 3 persistent-residual rules & $T{=}200$--$400$     & Progressive jumps, then locked \\
3D & diamoeba3d                  & $T{=}40$             & Initial transient, then locked \\
\bottomrule
\end{tabular}
\end{table}

In all tested cases, the NML-selected period makes finitely many transitions then stabilizes over the observed range, consistent with Theorem~\ref{thm:stabilization}.
The margins grow after stabilization in all cases, providing empirical evidence that the asymptotic regime has been reached, though finite data cannot prove this.
The stabilization rate depends on how quickly the asymptotic NLL rates $\lambda_c$ separate: 1D rules separate immediately, 3D rules after one correction, and 2D persistent-residual rules (where residuals have their own periodic structure) take longest.

%======================================================================
\section{Application: Persistent Structured Defects in 2D}
%======================================================================

\subsection{400-Step Verification}

We ran 400-step simulations ($64 \times 64$, seed 11) and measured defect counts in 100-step windows:

\begin{table}[h]
\centering
\caption{400-step defect persistence.}
\begin{tabular}{llcccc}
\toprule
Rule & Period & Window & Mean Defects & Std & Range \\
\midrule
S24/B11 & 2 & $t{=}50$--$100$  & 14.8 & 1.3 & [13, 17] \\
S24/B11 & 2 & $t{=}300$--$400$ & 14.8 & 1.3 & [13, 17] \\
S11/B37 & 4 & $t{=}50$--$100$  & 15.5 & 5.2 & [10, 31] \\
S11/B37 & 4 & $t{=}300$--$400$ & 12.6 & 2.0 & [10, 15] \\
S37/B11 & 2 & $t{=}50$--$100$  & 28.6 & 4.2 & [21, 37] \\
S37/B11 & 2 & $t{=}300$--$400$ & 28.1 & 3.9 & [18, 37] \\
\bottomrule
\end{tabular}
\end{table}

All three rules show stable defect populations from $t{=}100$ onward.

\subsection{Size Scaling}

5 seeds, grid sizes 32--192, 100 steps.

\paragraph{S37/B11---Extensive scaling.}

\begin{table}[h]
\centering
\caption{S37/B11 size scaling.}
\begin{tabular}{lccc}
\toprule
Grid & Mean Rate & Late Defects & Defect Density \\
\midrule
$32 \times 32$   & 3.47\% & $11.4 \pm 5.3$  & 0.011 \\
$64 \times 64$   & 3.26\% & $45.7 \pm 6.6$  & 0.011 \\
$96 \times 96$   & 3.11\% & $101.9 \pm 8.0$ & 0.011 \\
$128 \times 128$ & 3.33\% & $196.5 \pm 10.6$ & 0.012 \\
$192 \times 192$ & 3.14\% & $428.5 \pm 32.4$ & 0.012 \\
\bottomrule
\end{tabular}
\end{table}

Defect density is approximately constant (${\sim}0.011$) across 5 seeds and 5 grid sizes, consistent with a bulk (extensive) phenomenon.

\paragraph{S11/B37---Sub-extensive scaling.}

\begin{table}[h]
\centering
\caption{S11/B37 size scaling.}
\begin{tabular}{lccc}
\toprule
Grid & Mean Rate & Late Defects & Defect Density \\
\midrule
$32 \times 32$   & 1.46\% & $1.6 \pm 0.6$   & 0.0016 \\
$64 \times 64$   & 1.81\% & $11.1 \pm 6.3$  & 0.0027 \\
$96 \times 96$   & 1.71\% & $18.6 \pm 6.2$  & 0.0020 \\
$128 \times 128$ & 1.58\% & $32.7 \pm 4.1$  & 0.0020 \\
$192 \times 192$ & 1.69\% & $79.5 \pm 16.7$ & 0.0022 \\
\bottomrule
\end{tabular}
\end{table}

Defect density drops from $64 \times 64$ to $96 \times 96$ then stabilizes around 0.002, consistent with a sub-extensive population.

%======================================================================
\section{Discussion}
%======================================================================

\subsection{What the Theory Provides}

The orbit-class reduction (Theorem~\ref{thm:projection}) transforms a seemingly complex spatiotemporal fitting problem into standard Bernoulli estimation. This enables:

\begin{enumerate}
    \item \textbf{Exact optimal decomposition} in $O(n)$ time for any dimension.
    \item \textbf{A formal explanation} of overcapacity (Theorem~\ref{thm:monotonicity}), showing why naive defect-rate ranking is inadequate.
    \item \textbf{A stabilization theorem} (Theorem~\ref{thm:stabilization}) proving that NML model selection over the finite candidate set eventually locks to a unique winner, with a single assumption (ergodic orbit-class frequencies) and no dependence on defect geometry.
    \item \textbf{A nonidentifiability result} (Theorem~\ref{thm:nonident}) showing that background period recovery is impossible in general, with recovery guaranteed only when residuals are aperiodic (Theorem~\ref{thm:ident}).
\end{enumerate}

The key insight is that the \emph{same} theoretical framework---orbit classes $\to$ Bernoulli estimation $\to$ NML stabilization---applies without modification across spatial dimensions.

\subsection{NML vs Legacy MDL}

The previous formulation used a two-part MDL criterion: approximate parametric penalty $\frac{k}{2}\log_2(T/p)$ plus run-length codelength $L_{\RL}$ of the defect mask.
This had three weaknesses:

\begin{enumerate}
    \item \textbf{Model mismatch}: the RL codelength is not the data-fit term of the Bernoulli model implied by the orbit-class structure. NLL is.
    \item \textbf{Geometry dependence}: RL codelength depends on the spatial arrangement of defects and the flattening order, making it non-intrinsic.
    \item \textbf{Extra assumption}: the convergence proof required an $\Omega(T)$ runs condition to ensure $L_{\RL} = \Theta(T)$, which NML does not need.
\end{enumerate}

In practice, NML and MDL agree on period selection for most rules (e.g., both select period~4 for Rule~54).
They diverge for 2D rules with periodic defects, where NML selects lower, more parsimonious periods.

We retain RL and LZ4 codelengths as geometric diagnostics for characterizing defect spatial structure after model selection.

\subsection{What the Framework Does Not Provide}

The method identifies \emph{where} defects are but not \emph{how} they interact.
It does not recover domain/particle structure in the computational mechanics sense \cite{crutchfield1997,hanson1992}.
The fitted backgrounds are nearest periodic fields, not exact CA orbits.
Component-level lifetime and interaction analysis is future work.

\subsection{Limitations}

\begin{enumerate}
    \setcounter{enumi}{-1}
    \item \textbf{Monotonicity requires velocity matching.}
    Theorem~\ref{thm:monotonicity} holds along constant-velocity chains ($s_2 = m \cdot s_1 \bmod D$), not arbitrary period multiples with the same shift.
    For shift-zero scanning (the empirically dominant case), all divisibility chains satisfy the condition.
    For nonzero shifts, this restriction matters.

    \item \textbf{Shift scanning adds little for this rule family.}
    All 2D best fits have shift $(0,0)$.
    The theoretical framework supports nonzero shifts, but the empirical gain is rule-dependent.

    \item \textbf{NML selects the best-compressing period, not necessarily the ``true'' period.}
    When residuals have their own periodicity, NML absorbs them into higher-period templates (Theorem~\ref{thm:nonident}).
    This is principled but may not match a physicist's intuition.

    \item \textbf{Asymptotic NML complexity.}
    The $\frac{1}{2}\log_2 n_j$ complexity term is an asymptotic approximation to the Bernoulli NML normalizer.
    For very small orbit classes ($n_j < 10$), the approximation is less tight---relevant only at the smallest $T$ values.

    \item \textbf{Bernoulli i.i.d.\ assumption.}
    Within each orbit class, the NML criterion treats observations as i.i.d.\ Bernoulli.
    In reality, successive observations within an orbit class are separated by $p$ time steps and may have temporal correlations.
    Under ergodicity (A1), these correlations do not affect the asymptotic rate, but they may affect finite-sample behavior.
\end{enumerate}

\subsection{Future Work}

\begin{enumerate}
    \item \textbf{Component tracking}: lifetimes, velocities, collision catalogs.
    \item \textbf{Direct comparison} with local causal states \cite{rupe2018} on Rule~54 and 2D rules.
    \item \textbf{Larger surveys}: extend to 1000+ 2D rules and broader 3D rule families.
    \item \textbf{Finite-sample corrections}: exact (non-asymptotic) Bernoulli NML normalizer, or plug-in corrections for small orbit classes.
    \item \textbf{Shift optimization}: joint NML optimization over both shift and period.
\end{enumerate}

%======================================================================
\section{Conclusion}
%======================================================================

We have proved that Bernoulli NML model selection over relative-periodic CA backgrounds stabilizes: the selected period eventually locks to a unique winner as the observation window grows (Theorem~\ref{thm:stabilization}), though the selected period need not correspond to any external notion of ``true background period'' (Theorem~\ref{thm:nonident}).
The proof requires only ergodic orbit-class frequencies---no assumption on defect geometry---and is dimension-agnostic, resting on the orbit-class reduction that decomposes spatiotemporal fitting into independent Bernoulli estimation.
When residuals are aperiodic, NML recovers the true background period (Theorem~\ref{thm:ident}); when residuals are periodic, NML correctly absorbs them into the template.

Cross-dimensional experiments on 1D elementary CA, 2D totalistic rules, and 3D totalistic rules are consistent with the stabilization prediction, with NML-selected periods locking and margins growing monotonically over the observed range.
As an application, the framework identifies 2D rules with persistent structured projection residuals exhibiting extensive scaling---candidates for further computational-mechanics analysis.

%======================================================================
\appendix
%======================================================================

\section{Reproducibility}\label{app:repro}

All code is open-source. Key parameters:
\begin{itemize}
    \item 1D ECA: width${}=192$, steps${}=144$, density${}=0.5$, seed${}=11$, shifts${}=-6..{+}6$, periods${}=1..10$
    \item 1D convergence: $T \in \{50, 100, 200, 400, 600, 800\}$, periods${}=1..10$
    \item 2D survey: $48 \times 48$, steps${}=40$, 1{,}050 candidates (773 non-trivial), range width $\leq 4$
    \item 2D convergence: $64 \times 64$, $T \in \{60, 100, 200, 400, 600, 800\}$, periods${}=1..8$
    \item 3D convergence: $16^3$, $T \in \{10, 20, 40, 60, 80\}$, periods${}=1..6$
    \item Multi-seed: seeds $= \{11, 42, 73, 99, 137, 200, 314, 500, 777, 1024\}$
    \item Size scaling: grids $\in \{32, 64, 96, 128, 192\}$, steps${}=100$, 5 seeds
    \item 400-step verification: $64 \times 64$, steps${}=400$, seed${}=11$
    \item Model selector: Bernoulli NML (orbit-class NLL + asymptotic parametric complexity)
\end{itemize}

\section{Geometric Diagnostics}\label{app:geom}

Run-length and LZ4 codelengths are used as geometric diagnostics, not for model selection.
Example: two synthetic masks of size 512, each with 64 defects---clustered (42 RL bits) vs scattered (251 RL bits), a $6\times$ ratio at identical Hamming weight.
This illustrates Proposition~\ref{prop:rl-separation} and motivates RL/LZ4 as post-selection defect geometry characterizers.

\section{NML vs Legacy MDL Period Selection}\label{app:nml-mdl}

\begin{table}[h]
\centering
\caption{NML vs MDL period selection comparison.}
\begin{tabular}{llll}
\toprule
Rule & MDL Period & NML Period & Agreement \\
\midrule
ECA-54 (1D)          & 4 & 4 & Yes \\
ECA-110 (1D)         & 7 & 7 & Yes \\
ECA-30 (1D)          & 1 & 1 & Yes \\
S24/B11 (2D, $T{=}800$) & 6 & 2 & No---NML more parsimonious \\
S11/B37 (2D, $T{=}800$) & 8 & 4 & No---NML more parsimonious \\
S37/B11 (2D, $T{=}800$) & 6 & 2 & No---NML more parsimonious \\
diamoeba3d (3D)      & 1 & 1 & Yes \\
\bottomrule
\end{tabular}
\end{table}

NML and MDL agree for 1D and 3D rules but diverge for 2D rules with periodic defects, where the RL data-fit artifact in MDL inflates scores for higher periods.

\section{RL Rate Convergence for Finite Deterministic CA}\label{app:rl}

This appendix closes the convergence gap noted in Remark~\ref{rem:rl} by proving, from first principles, that the per-site run-length rate converges for finite deterministic CA.

\begin{lemma}[RL convergence for eventually periodic sequences]\label{lem:rl-periodic}
Let $b_1, b_2, \ldots$ be a binary sequence that is eventually periodic: there exist $T_0 \geq 0$ and $Q \geq 1$ such that $b_{n+Q} = b_n$ for all $n > T_0$.
Then
\[
\lim_{N \to \infty} \frac{L_{\RL}(b_1 \cdots b_N)}{N} = \frac{C}{Q},
\]
where $C$ is the per-period RL cost of the periodic part, defined as follows.
Let the period pattern $b_{T_0+1}, \ldots, b_{T_0+Q}$ have runs $r_1, \ldots, r_R$ (under Elias gamma coding $\gamma$).
\begin{itemize}
    \item If $b_{T_0+Q} \neq b_{T_0+1}$ (no seam merging): $C = \sum_{i=1}^{R} \gamma(r_i)$.
    \item If $b_{T_0+Q} = b_{T_0+1}$ (seam merging): $C = \gamma(r_R + r_1) + \sum_{i=2}^{R-1} \gamma(r_i)$.
\end{itemize}
\end{lemma}

\begin{proof}
Write $N = T_0 + M Q + R'$ where $0 \leq R' < Q$ and $M = \lfloor (N - T_0)/Q \rfloor$.
The RL codelength decomposes as
\[
L_{\RL}(b_1 \cdots b_N) = L_{\text{transient}} + M \cdot C + O(\log Q),
\]
where $L_{\text{transient}}$ accounts for the prefix $b_1 \cdots b_{T_0}$ and any boundary run merging with the periodic part (a fixed $O(1)$ cost), and the $O(\log Q)$ term accounts for a potentially truncated final period.
Therefore
\[
\frac{L_{\RL}(N)}{N} = \frac{M \cdot C + O(1)}{T_0 + MQ + R'} \to \frac{C}{Q}
\]
as $M \to \infty$.
\end{proof}

\begin{theorem}[RL Rate Convergence for Finite CA]\label{thm:rl-convergence}
Let $U$ be the spacetime of a deterministic CA on a finite spatial domain $S$ with $|S| = D$.
Let $(p, s)$ be a candidate model.
Assume:
\begin{itemize}
    \item[\textup{(NT)}] \emph{No exact ties}: for each orbit class $j$, the limiting empirical frequency $\theta_j^* \neq 1/2$.
\end{itemize}
Then the per-site RL rate converges:
\[
\lim_{T \to \infty} \frac{L_{\RL}(M^*_c(T))}{T \cdot D} \quad \text{exists}.
\]
\end{theorem}

\begin{proof}
The proof chains four facts.

\medskip\noindent
\textbf{Step 1 (Eventual periodicity of the CA).}
The CA operates on the finite state space $\{0,1\}^D$.
By the pigeonhole principle, the state sequence $U[0], U[1], \ldots$ is eventually periodic: there exist $T_1$ and $P$ such that $U[t + P] = U[t]$ for all $t \geq T_1$.

\medskip\noindent
\textbf{Step 2 (Majority vote stabilization).}
For each orbit class $j$, the empirical frequency $\hat\theta_j(T)$ converges to $\theta_j^*$ as $T \to \infty$ (because the per-period contribution is fixed after the transient, and the transient contribution is $O(1/T)$).
By assumption (NT), $\theta_j^* \neq 1/2$, so the majority vote $b_j^* = [\theta_j^* > 1/2]$ is eventually stable: there exists $T_2$ such that the majority vote is constant for all $T \geq T_2$.

\medskip\noindent
\textbf{Step 3 (Eventual periodicity of the defect mask).}
Set $T_0 = \max(T_1, T_2)$.
For $t \geq T_0$:
\begin{itemize}
    \item The CA state $U[t]$ is periodic with period $P$.
    \item The background $B^*(T)$ is fixed (for $T \geq T_2$), so $M[t, x] = U[t, x] \oplus B^*[t, x]$ is determined by $U[t, x]$ and the fixed template.
    \item The template value $B^*[t, x]$ depends on $(t \bmod p,\; (x - \lfloor t/p \rfloor \cdot s) \bmod D)$, which is periodic in $t$ with period $\Lambda_B = p \cdot D / \gcd(|s|, D)$.
    \item The mask row $M[t, :]$ is periodic with period $\Lambda = \operatorname{lcm}(P, \Lambda_B)$ for $t \geq T_0$.
    \item The flattened mask is eventually periodic with period $\Lambda \cdot D$.
\end{itemize}

\medskip\noindent
\textbf{Step 4 (RL convergence).}
By Lemma~\ref{lem:rl-periodic}, $L_{\RL}(\text{flattened mask}) / N$ converges as $N = T \cdot D \to \infty$.
\end{proof}

\begin{corollary}[Unconditional NML convergence]\label{cor:nml-unconditional}
For a deterministic CA on a finite spatial domain, the per-site NLL converges \emph{unconditionally} (no assumption \textup{(NT)} needed):
\[
\lim_{T \to \infty} \frac{\NLL_c(T)}{T \cdot D} = h_c \quad \text{exists}.
\]
\end{corollary}

\begin{proof}
The per-orbit-class frequency $\hat\theta_j(T) \to \theta_j^*$ follows from Steps~1 and~2 of Theorem~\ref{thm:rl-convergence}, without requiring (NT).
Binary entropy $\Hb$ is continuous, so $n_j \cdot \Hb(\hat\theta_j(T)) / (T \cdot D) \to (1/p) \cdot \Hb(\theta_j^*)$.
Summing over $j$ gives the limit.

When $\theta_j^* = 1/2$, the NLL contribution is $n_j \cdot \Hb(1/2) = n_j$ bits, regardless of the tie-breaking direction.
So the NLL does not depend on (NT).
Combined with the $O(\log T)$ complexity penalty, $\NML(c, T) / (T \cdot D) \to h_c$, and by Theorem~\ref{thm:stabilization}, model selection stabilizes.
No assumption beyond finiteness of the CA is needed.
\end{proof}

\begin{remark}[When \textup{(NT)} fails: RL may not converge]\label{rem:nt-failure}
When an orbit class $j$ has $\theta_j^* = 1/2$ exactly, the majority vote may oscillate: as $T$ increases, the cumulative count can cross the $0.5$ threshold back and forth.
Each flip of the majority vote changes the mask at all $O(T)$ sites in orbit class $j$, potentially changing $L_{\RL}$ by $\Theta(T)$.
So $L_{\RL}/N$ may fail to converge.

The exact-tie condition is \emph{non-generic} (it requires the count of 1s in an orbit class over one CA period to be exactly half the class size) but is not impossible---Rule~30 on width~8 has all classes with $\theta^* = 0.5$ exactly.

This is another advantage of the NML score over the RL score: NML converges unconditionally, while RL requires the no-ties assumption.
\end{remark}

%======================================================================
\begin{thebibliography}{10}

\bibitem{crutchfield1997}
J.\,P.~Crutchfield and J.\,E.~Hanson,
``Computational mechanics of cellular automata: An example,''
\emph{Physica D}, vol.~103, pp.~169--189, 1997.

\bibitem{hanson1992}
J.\,E.~Hanson and J.\,P.~Crutchfield,
``The attractor-basin portrait of a cellular automaton,''
\emph{J.\ Statistical Physics}, vol.~66, pp.~1415--1462, 1992.

\bibitem{redeker2017}
M.~Redeker,
``A language for particle interactions in Rule~54 and other cellular automata,''
\emph{Complex Systems}, vol.~26, no.~1, pp.~1--32, 2017.

\bibitem{rupe2018}
A.~Rupe and J.\,P.~Crutchfield,
``Local causal states and discrete coherent structures,''
\emph{Chaos}, vol.~28, 075312, 2018.

\bibitem{packard1985}
N.\,H.~Packard and S.~Wolfram,
``Two-dimensional cellular automata,''
\emph{J.\ Statistical Physics}, vol.~38, pp.~901--946, 1985.

\bibitem{boccara1999}
N.~Boccara and M.~Roger,
``Totalistic two-dimensional cellular automata exhibiting global periodic behavior,''
\emph{Int.\ J.\ Modern Physics~C}, vol.~10, no.~6, pp.~1051--1066, 1999.

\bibitem{shalizi2006}
C.\,R.~Shalizi, R.~Haslinger, J.-B.~Rouquier, K.\,L.~Klinkner, and C.~Moore,
``Automatic filters for the detection of coherent structure in spatiotemporal systems,''
\emph{Physical Review~E}, vol.~73, 036104, 2006.

\bibitem{zenil2010}
H.~Zenil,
``Compression-based investigation of the dynamical properties of cellular automata and other systems,''
\emph{Complex Systems}, vol.~19, no.~1, pp.~1--28, 2010.

\bibitem{grunwald2007}
P.~Gr\"unwald,
\emph{The Minimum Description Length Principle},
MIT Press, 2007.

\bibitem{shtarkov1987}
Y.~Shtarkov,
``Universal sequential coding of single messages,''
\emph{Problems of Information Transmission}, vol.~23, pp.~3--17, 1987.

\end{thebibliography}

\end{document}
